%=============================================================================
% SECTION 10: CAVITY-ATOM LPI TEST
% Part VI: Laboratory Tests
% Equations: (54)-(56)
% Estimated pages: 6-8
%=============================================================================

\section{Cavity-Atom Redshift Tests}\label{sec:cavity}

The cavity--atom comparison remains part of the DFD laboratory program, but its role changed substantially once the optical-metric constitutive chain was treated consistently.  Earlier internal drafts effectively slowed light while holding the cavity spacer fixed, producing an order-unity LPI slope.  That is not the correct DFD calculation.  In the corrected treatment, the same optical metric that changes photon propagation also changes Coulomb binding, lattice spacing, and hence the cavity length.  The leading geometric response of cavity and atomic sectors cancels at tree level.

What survives is a residual, screened signal.  This makes the cavity--atom channel \emph{harder} as an experiment but also cleaner as a precision residual test.

%-----------------------------------------------------------------------------
\subsection{Formal Constitutive Proof of the Cancellation}\label{sec:cavity-proof}
%-----------------------------------------------------------------------------

The cancellation can be organized as a short formal derivation.

\paragraph{Step 1: optical metric and constitutive relations.}
DFD posits the optical metric
\begin{equation}\label{eq:optical-metric-cavity}
d\tilde s^2=-\frac{c^2}{n^2}dt^2+d\mathbf{x}^2,
\qquad n=e^{\psi}.
\end{equation}
Through the Tamm--Plebanski construction, this metric defines effective vacuum constitutive relations
\begin{equation}\label{eq:tamm-plebanski}
\varepsilon_{\rm eff}=\varepsilon_0 e^{+\psi},
\qquad
\mu_{\rm eff}=\mu_0 e^{+\psi}.
\end{equation}
The medium is impedance-matched, and the local phase velocity is
\begin{equation}
v_{\rm ph}=\frac{1}{\sqrt{\varepsilon_{\rm eff}\mu_{\rm eff}}}=ce^{-\psi}.
\end{equation}

\paragraph{Step 2: Coulomb binding changes with the same constitutive chain.}
Virtual photons feel the same optical medium, so the static Coulomb potential scales as
\begin{equation}\label{eq:coulomb-psi}
V(r)=\frac{e^2}{4\pi\varepsilon_{\rm eff}r}
=\frac{e^2}{4\pi\varepsilon_0 r}e^{-\psi}.
\end{equation}
The local fine-structure constant at tree level is therefore unchanged:
\begin{equation}\label{eq:alpha-tree-cancel}
\alpha(\psi)=\frac{e^2}{4\pi\varepsilon_{\rm eff}\hbar c_{\rm local}}
=\alpha_0,
\end{equation}
because the factors from $\varepsilon_{\rm eff}$ and $c_{\rm local}=ce^{-\psi}$ cancel.

\paragraph{Step 3: the atomic length scale expands.}
With $\alpha$ unchanged at tree level, the Bohr radius scales as
\begin{equation}\label{eq:bohr-psi}
a_0(\psi)=\frac{\hbar}{m_ec_{\rm local}\alpha}=a_0^{(0)}e^{+\psi}.
\end{equation}
Thus the microscopic electromagnetic length scale expands in stronger field.

\paragraph{Step 4: the cavity length follows the same electromagnetic scale.}
A Fabry--P\'erot cavity resonance obeys
\begin{equation}
f_{\rm cav}=\frac{mc_{\rm local}}{2L(\psi)}.
\end{equation}
For an electromagnetic solid spacer, the lattice constant and therefore $L$ scale with the Bohr radius, so $L\propto e^{+\psi}$ while $c_{\rm local}\propto e^{-\psi}$.  Hence
\begin{equation}\label{eq:fcav-psi}
f_{\rm cav}\propto \frac{e^{-\psi}}{e^{+\psi}}=e^{-2\psi}.
\end{equation}
Atomic transition frequencies scale with the same leading factor, $f_{\rm atom}\propto e^{-2\psi}$, up to channel-dependent residual sensitivities.

\paragraph{Convention note.}
The $e^{-2\psi}$ scaling above is in coordinate time, derived from the optical-metric constitutive chain ($c_{\rm local}\propto e^{-\psi}$, Bohr radius $\propto e^{+\psi}$, $E_n \propto m_e c_{\rm local}^2\alpha^2 \propto e^{-2\psi}$).  Section~\ref{sec:ppn-redshift} quotes $\nu\propto e^{-\psi/2}$, which is the gravitational redshift factor from the physical metric $g_{00}=-e^{-\psi}$.  The key point is not that the individual exponents match---they refer to different quantities---but that the \emph{same} universal coordinate-to-proper conversion multiplies both cavity and atomic frequencies.  Therefore the ratio $R = f_{\rm cav}/f_{\rm atom}$ is convention-independent, and the tree-level cancellation holds regardless of which clock convention is adopted.

\paragraph{Tree-level result.}
The leading geometric ratio is therefore constant:
\begin{equation}\label{eq:ratio-cancel}
\boxed{\;R\equiv \frac{f_{\rm cav}}{f_{\rm atom}}=\text{const. at tree level},\qquad \xi_{\rm geom}=0\;}
\end{equation}
The universal geometric redshift cancels.  Any surviving cavity--atom signal must come from a \emph{residual} channel, not from an order-unity tree-level effect.

%-----------------------------------------------------------------------------
\subsection{What Survives Physically}\label{sec:cavity-residual}
%-----------------------------------------------------------------------------

Once the tree-level cancellation is enforced, the cavity--atom observable is naturally written as
\begin{equation}\label{eq:cavity-atom-slope}
\frac{\Delta R}{R}=\xi_{\rm LPI}^{\rm res}\frac{\Delta\Phi}{c^2},
\end{equation}
with
\begin{equation}\label{eq:xi-prediction}
\boxed{\xi_{\rm LPI}^{\rm GR}=0,\quad \xi_{\rm LPI}^{\rm DFD}=\text{screened residual}.}
\end{equation}
For ordinary terrestrial experiments this residual is small because the local environment sits deep in the screened regime.

\paragraph{Interpretation.}
This is not a failure of the master program; it is a correction of the measurement channel.  The cavity--atom comparison remains valuable precisely because it can isolate a non-metric residual if the sensitivity frontier is pushed far enough.

%-----------------------------------------------------------------------------
\subsection{Three Independent Empirical Checks}\label{sec:cavity-checks}
%-----------------------------------------------------------------------------

The geometric-cancellation picture is not just a pretty derivation.  Three independent data streams push in the same direction.

\paragraph{Check 1: fine-structure splitting.}
If the geometric unscreened picture were right, the ratio of two transitions with different $\alpha$ sensitivities inside the same atom would show an annual modulation of order $\Delta S^\alpha\,\delta\psi_{\rm annual}\sim 10^{-10}$.  Precision spectroscopy constrains such effects at the $\lesssim10^{-17}$ level.  The naive unscreened geometric scenario is therefore ruled out by more than seven orders of magnitude.

\paragraph{Check 2: PTB Yb$^+$ E3/E2.}
The same-ion E3/E2 comparison~\cite{Lange2021} is exactly the sort of experiment that would have seen the old unscreened cavity-style logic if it were real.  Instead, the observed result is null at a level that rules out the naive geometric expectation by roughly two orders of magnitude and forces any viable theory into a much smaller residual regime.

\paragraph{Check 3: BACON optical network.}
The BACON collaboration measured Al$^+$/Sr/Yb frequency ratios with uncertainties at or below $8\times10^{-18}$~\cite{Beloy2021}.  A naive geometric annual signal in Yb/Sr would be of order $4\times10^{-11}$, absurdly larger than the observed stability.  This is effectively a million-fold exclusion of the unscreened geometric cavity/atom picture.

These three checks all point the same way: the order-unity tree-level picture is dead; only a screened residual can survive.

%-----------------------------------------------------------------------------
\subsection{BACON and the Screening Regime}\label{sec:cavity-bacon}
%-----------------------------------------------------------------------------

BACON does more than kill the naive tree-level picture.  It also constrains \emph{how} screening should be evaluated.

\paragraph{Solar-orbit screening fails.}
If one evaluates the residual coupling at the solar-orbit acceleration, Eq.~\eqref{eq:kalpha-eff} gives roughly
\begin{equation}
k_\alpha^{\rm eff}(1\,\mathrm{AU})\approx 2.4\times10^{-5}.
\end{equation}
For Yb/Sr, with $\Delta S^\alpha\approx 0.25$, the implied annual signal is then
\begin{equation}
\left(\frac{\delta R}{R}\right)_{\rm Yb/Sr}\approx 0.25\times 2.4\times10^{-5}\times1.65\times10^{-10}
\approx 10^{-15}.
\end{equation}
BACON's weighted scatter for Yb/Sr is about $1.1\times10^{-17}$, so this solar-orbit-screened scenario is excluded by roughly two orders of magnitude.

\paragraph{Earth-surface screening survives.}
If instead the local gravitational environment controls the screening, then at Earth's surface
\begin{equation}
k_\alpha^{\rm eff}(\oplus)\approx 6\times10^{-7},
\end{equation}
and the same Yb/Sr estimate becomes
\begin{equation}
\left(\frac{\delta R}{R}\right)_{\rm Yb/Sr}\approx 2.5\times10^{-17},
\end{equation}
which is comparable to the observed between-day variability and therefore not excluded by BACON.

\paragraph{Operational conclusion.}
The residual screening must be evaluated using the \emph{local} background acceleration, in agreement with the screening analysis built from the BACON network.  This is a nontrivial quantitative result and should be regarded as one of the main takeaways of the corrected cavity--atom program.

%-----------------------------------------------------------------------------
\subsection{Sector-Resolved Parameterization}\label{sec:cavity-sectors}
%-----------------------------------------------------------------------------

The cavity--atom channel still benefits from a sectoral bookkeeping language.  Write
\begin{align}
\left(\frac{\Delta f}{f}\right)_{\rm cav}^{(M)} &= (\alpha_{\rm w}-\alpha_L^{(M)})\frac{\Delta\Phi}{c^2}, \\
\left(\frac{\Delta f}{f}\right)_{\rm atom}^{(S)} &= \alpha_{\rm atom}^{(S)}\frac{\Delta\Phi}{c^2}.
\end{align}
Only differences are observable.  With two cavity materials (for example ULE and Si) and two atomic species (for example Sr and Yb), the directly identifiable combinations are
\begin{align}
\delta_{\rm tot} &\equiv \alpha_{\rm w}-\alpha_L^{\rm ULE}-\alpha_{\rm atom}^{\rm Sr}, \\
\delta_L &\equiv \alpha_L^{\rm Si}-\alpha_L^{\rm ULE}, \\
\delta_{\rm atom} &\equiv \alpha_{\rm atom}^{\rm Yb}-\alpha_{\rm atom}^{\rm Sr}.
\end{align}
This remains useful for a future high-precision residual measurement even after the tree-level cancellation is imposed.

%-----------------------------------------------------------------------------
\subsection{The 4$\to$3 GLS Protocol}\label{sec:cavity-gls}
%-----------------------------------------------------------------------------

The four basic cavity--atom slopes still map cleanly onto three independent sector combinations:
\begin{table}[h]
\centering
\caption{Mapping of measured cavity--atom ratios to sector parameters.}
\label{tab:gls-mapping}
\begin{ruledtabular}
\begin{tabular}{lcc}
Measured slope & Combination & Parameter \\
\hline
ULE/Sr & $\alpha_{\rm w}-\alpha_L^{\rm ULE}-\alpha_{\rm atom}^{\rm Sr}$ & $\delta_{\rm tot}$ \\
Si/Sr & $\alpha_{\rm w}-\alpha_L^{\rm Si}-\alpha_{\rm atom}^{\rm Sr}$ & $\delta_{\rm tot}+\delta_L$ \\
ULE/Yb & $\alpha_{\rm w}-\alpha_L^{\rm ULE}-\alpha_{\rm atom}^{\rm Yb}$ & $\delta_{\rm tot}+\delta_{\rm atom}$ \\
Si/Yb & $\alpha_{\rm w}-\alpha_L^{\rm Si}-\alpha_{\rm atom}^{\rm Yb}$ & $\delta_{\rm tot}+\delta_L+\delta_{\rm atom}$ \\
\end{tabular}
\end{ruledtabular}
\end{table}

The redundancy provides a built-in closure relation and remains valuable even though the target signal is now residual rather than order unity.

%-----------------------------------------------------------------------------
\subsection{Experimental Concept and Controls}\label{sec:cavity-experiment}
%-----------------------------------------------------------------------------

The experimental architecture developed in earlier drafts still has value and is retained here because the correction changed the \emph{amplitude}, not the measurement logic.

\paragraph{Hardware.}
\begin{itemize}
\item two evacuated optical cavities (for example ULE and cryogenic Si) with PDH-locked lasers;
\item co-located Sr and Yb optical lattice clocks;
\item a self-referenced frequency comb measuring all four ratios simultaneously;
\item vertical relocation or a dual-station geometry providing a known potential difference.
\end{itemize}

\paragraph{Dispersion control.}
The dual-wavelength check remains essential.  DFD's optical metric is nondispersive in the minimal formulation, so any large wavelength dependence would diagnose ordinary optical systematics rather than a gravitational effect.  Causality constrains material dispersion via the Kramers--Kronig relation:
\begin{equation}\label{eq:KK-bound}
\left|\frac{\partial\ln n}{\partial\ln\omega}\right|
 \lesssim \frac{2}{\pi}\frac{\omega}{\Omega}
           \frac{\alpha_0 L_{\rm mat}}{\mathcal{F}},
\end{equation}
where $\mathcal{F}$ is the cavity finesse, $L_{\rm mat}$ the material path length, $\alpha_0$ the absorption coefficient, and $\Omega$ the detuning to the nearest material resonance.  For crystalline mirror coatings and ULE glass near optical-clock frequencies ($\alpha_0 < 10^{-4}$, $\Omega/\omega > 10^{-2}$), this yields $|\xi - 1| < 10^{-8}$---far below experimental reach.

\paragraph{Cavity mechanics.}
Vertical transport changes gravitational loading on the cavity spacer.  Controls include: elastic modeling to null first-order sag; 180$^\circ$ orientation flips at each height (mechanical artifacts change sign, gravitational effects do not); and a platform tilt budget maintained at ${<}100\,\mu$rad.  Gravitational sag contributes $\alpha_{\rm grav} \sim 10^{-9}$ for ULE, elastic coupling ${<}10^{-14}$ for $10^{-6}g$ perturbations, and thermoelastic drift cancels in common-mode ratios.  The combined effective length-change bound is $|\alpha_L^M| \lesssim 10^{-8}$.

\paragraph{Environmental and noise budget.}
Temperature stability ${<}10$ mK, pressure ${<}10^{-2}$ mbar, magnetic field drift ${<}10\,\mu$T with periodic reversal.  The ratio Allan variance is modeled as $\sigma_y^2(\tau) = h_{-1}/\tau + h_0 + h_1\tau$, with typical values: white frequency $h_{-1} \sim 10^{-32}$ (300\,s windows), flicker $h_0 \sim 10^{-34}$, random walk $h_1 \sim 10^{-38}$.  The dominant term is white noise.

\paragraph{Thermal rejection.}
Silicon cavities have $dn/dT \sim 10^{-4}$/K; with $\delta T < 10$ mK the fractional contribution is ${<}10^{-6}$.  ULE has CTE $\approx 0$ near 30$^\circ$C; silicon near 124 K has CTE $\approx 0$.  Operating at these zero-crossings suppresses length changes.  Any residual dispersion from coating thermal effects appears differently at two wavelengths, bounding the dispersion systematic to $|\epsilon_{\rm disp}| \lesssim 10\%$.  Total thermal target: ${<}3 \times 10^{-16}$, achievable with demonstrated technology.

%-----------------------------------------------------------------------------
\subsection{Expected Signal and Sensitivity}\label{sec:cavity-sensitivity}
%-----------------------------------------------------------------------------

For a height separation $\Delta h$,
\begin{equation}\label{eq:cavity-signal}
\frac{g\Delta h}{c^2}\approx 1.1\times10^{-14}\left(\frac{\Delta h}{100\,\mathrm{m}}\right).
\end{equation}
After geometric cancellation, the cavity--atom observable inherits this factor \emph{and} a screened residual coefficient.  The terrestrial height-separated signal is therefore extremely small.

\paragraph{Consequence.}
The practical ranking of experiments changes:
\begin{itemize}
\item a terrestrial height-separated cavity--atom test is no longer a quick binary discriminator;
\item it becomes a demanding precision residual experiment, likely better matched to future long-baseline or space-based platforms;
\item multi-species clock and nuclear-clock programs move ahead of it in near-term priority.
\end{itemize}

%-----------------------------------------------------------------------------
\subsection{Current Status and Revised Priority}\label{sec:cavity-status}
%-----------------------------------------------------------------------------

No existing experiment has yet performed the full sector-resolved cavity--atom residual test at the required precision.  The correction therefore does \emph{not} mean the channel has been experimentally exhausted; it means the target has moved from ``large and immediate'' to ``clean but extremely small.''

\begin{tcolorbox}[colback=blue!5!white,colframe=blue!60!black,title=Revised Cavity--Atom Priority]
\textbf{Old picture:} first-line binary discriminator with $\xi_{\rm LPI}\sim1$.

\textbf{Corrected picture:} tree-level geometric cancellation; residual screened signal only.

\textbf{Revised ranking:}
\begin{enumerate}
\item Th-229/Sr and related nuclear-clock analyses
\item cross-species atomic clock comparisons
\item same-ion null checks that bound the pure $\alpha$ sector
\item height-separated cavity--atom residual tests
\end{enumerate}
\end{tcolorbox}

%-----------------------------------------------------------------------------
\subsection{Summary: Cavity--Atom as a Precision Residual Test}\label{sec:cavity-summary}
%-----------------------------------------------------------------------------

The corrected cavity--atom picture is now simple to state:
\begin{enumerate}
\item the optical metric implies constitutive relations through Tamm--Plebanski,
\item those constitutive relations alter both light propagation and electromagnetic binding,
\item the cavity spacer length therefore changes together with the local light speed,
\item the leading geometric cavity/atom response cancels at tree level,
\item only a residual screened signal survives.
\end{enumerate}

This section therefore remains in the master corpus for an important reason: it archives the complete logic of a channel that was once overstated and is now properly understood.  That makes the theory stronger, not weaker.
